% ============================================================
% Script de soutenance — Sheryne (diapos 3, 6 à 9)
% VERSION ADAPTÉE DYSLEXIE :
%   - police OpenDyslexic (fichiers dans fonts/), corps 12 pt
%   - interligne 1,6 et paragraphes espacés
%   - lignes courtes (marges larges), texte non justifié
%   - pas d'italique ; phrases courtes ; listes plutôt que pavés
% Compilation : xelatex script_sheryne_dyslexie.tex
% (xelatex obligatoire : la police est chargée via fontspec)
% ============================================================
\documentclass[12pt,a4paper]{article}

\usepackage[french]{babel}
\usepackage[margin=3.2cm]{geometry}
\usepackage{amsmath,amssymb}
\usepackage{booktabs}
\usepackage{parskip}
\usepackage{setspace}
\usepackage{fontspec}
\setmainfont{OpenDyslexic}[
Path           = fonts/,
UprightFont    = OpenDyslexic-Regular.otf,
BoldFont       = OpenDyslexic-Bold.otf,
ItalicFont     = OpenDyslexic-Italic.otf,
BoldItalicFont = OpenDyslexic-BoldItalic.otf,
]
\usepackage{enumitem}

\setstretch{1.6}
\raggedright
\setlength{\parskip}{0.9em}
\setlist{itemsep=0.5em, topsep=0.5em}

\newcommand{\diapo}[2]{\section*{Diapo #1 --- #2}}
\newenvironment{comprendre}%
{\par\medskip\noindent\hrulefill\par\nopagebreak\textbf{Pour bien comprendre :}\par}%
{\par\nopagebreak\noindent\hrulefill\par\medskip}

\title{\textbf{Script de soutenance --- Sheryne}\\
	\large Effet tunnel et temps de traversée d'une barrière\\
	\normalsize Diapos 3, 6--9 --- environ 5 minutes}
\author{}
\date{Soutenance juin 2026 --- oraux du 22 au 26 juin}

\begin{document}
	\maketitle
	
	\subsection*{Qui dit quoi}
	
	\begin{itemize}
		\item Myriam : diapos 1, 2, 4, 5. Le problème, le paquet d'ondes, la barrière. Environ 5 min.
		\item \textbf{Sheryne : diapos 3, 6 à 9. La motivation, la méthode numérique et les premiers résultats. Environ 5 min.}
		\item Mattéo : diapos 10 à 13. L'effet Hartman et la conclusion. Environ 5,5 min.
	\end{itemize}
	
	Total : environ 15,5 minutes, puis 10 minutes de questions.
	
	Les phrases \textbf{en gras} sont les messages essentiels.
	Si le temps presse, on peut raccourcir le reste, mais jamais le gras.
	
	% ------------------------------------------------------------
	\clearpage
	\diapo{3}{Motivation : de l'onde plane au paquet d'ondes (30 secondes)}
	
	\textbf{Texte à dire :}
	
	Merci Myriam.
	
	Avant le paquet gaussien, un mot sur son origine.
	
	Une onde plane seule n'est pas physique.
	Son module au carré est constant partout.
	\textbf{Elle n'est pas normalisable.}
	
	Pour une particule localisée, on superpose plusieurs ondes planes.
	Les interférences créent une enveloppe.
	\textbf{Cette enveloppe localise la particule.}
	C'est le battement que vous voyez sur la figure.
	
	On passe ensuite à une somme continue, sur tous les $k$.
	On obtient le paquet d'ondes gaussien.
	
	Je rends la parole à Myriam pour le détailler.
	
	% ------------------------------------------------------------
	\clearpage
	\diapo{6}{La méthode numérique : Crank--Nicolson (1,5 minute)}
	
	\textbf{Texte à dire :}
	
	Merci Myriam, je reprends.
	
	La théorie décrit un paquet libre, ou une onde plane face à la barrière.
	Mais un paquet complet face à une barrière : il n'y a pas de formule exacte.
	\textbf{Il faut donc résoudre l'équation de Schrödinger sur ordinateur.}
	
	Nous avons choisi le schéma de Crank--Nicolson.
	Son idée : \textbf{on prend la moyenne entre l'instant présent et l'instant suivant.}
	Cela donne l'équation à l'écran.
	Le coefficient $r$ contient les pas de temps et d'espace.
	Le coefficient $s_i$ contient le potentiel au point $x_i$.
	Le potentiel change selon la position, donc la diagonale du système change aussi.
	
	\textbf{Le code à droite fait exactement ça.}
	La ligne \texttt{s\_int} porte cette variation.
	C'est elle qui rend la diagonale du système différente à chaque point.
	
	Pourquoi ce schéma ? Trois raisons :
	
	\begin{itemize}
		\item \textbf{Il est unitaire.} La probabilité totale reste égale à 1.
		Vérifié à $10^{-6}$ près sur toute la simulation.
		Pour la mécanique quantique, c'est la propriété la plus importante.
		\item \textbf{Il est toujours stable}, quel que soit le pas de temps.
		Et il est précis à l'ordre 2 en temps et en espace.
		\item \textbf{Il est rapide.} Le système est tridiagonal.
		La méthode de Thomas le résout en un temps proportionnel au nombre de points.
		Vous la voyez appelée tout en bas du code.
	\end{itemize}
	
	Nos paramètres :
	
	\begin{itemize}
		\item grille de 1500 points sur $[-30, 30]$, 3000 pas de temps jusqu'à $t = 6$ ;
		\item paquet : $k_0 = 5$, donc énergie $E = 12{,}5$, départ en $x_0 = -10$ ;
		\item barrière : largeur $a = 1$, hauteur $V_0 = 15$.
	\end{itemize}
	
	\textbf{La barrière est plus haute que l'énergie du paquet : on est en régime tunnel.}
	
	\begin{comprendre}
		« Unitaire » : le schéma conserve la norme, comme la vraie évolution quantique.
		
		Le schéma d'Euler explicite, lui, augmente la norme à chaque pas.
		La probabilité dépasse 1 et la simulation explose.
		C'est la réponse si le jury demande : « pourquoi pas Euler ? »
		
		« Tridiagonal » : chaque point n'est couplé qu'à ses deux voisins.
		Cela vient de la dérivée seconde discrétisée.
	\end{comprendre}
	
	% ------------------------------------------------------------
	\clearpage
	\diapo{7}{Validation : particule libre (1 minute)}
	
	\textbf{Texte à dire :}
	
	Avant d'allumer la barrière, on teste le solveur sur un cas connu :
	la particule libre, avec $V_0 = 0$.
	
	Notre définition du temps de traversée :
	\textbf{on suit le pic de $|\psi|^2$.}
	On note l'instant où il passe en $x = 0$.
	Puis l'instant où il passe en $x = a$.
	La différence donne le temps numérique $\tau_0$.
	
	\textbf{C'est ce que fait la fonction à l'écran, \texttt{temps\_pic\_en}.}
	Elle repère où l'amplitude est maximale.
	Et renvoie l'instant correspondant.
	On l'appelle une fois à l'entrée, une fois à la sortie.
	
	La théorie est simple : distance divisée par vitesse.
	Donc $a / v_g = 0{,}2$.
	
	Le tableau montre l'accord :
	
	\begin{itemize}
		\item entrée mesurée à $t = 1{,}9987$, attendu : $2$ ;
		\item sortie mesurée à $t = 2{,}1907$, attendu : $2{,}2$ ;
		\item \textbf{donc $\tau_0 = 0{,}192$ contre $0{,}200$ : un écart de 4\,\%.}
	\end{itemize}
	
	Cet écart vient de la résolution en temps de la grille,
	et de l'étalement du paquet pendant le vol.
	La norme reste égale à 1 sur les 3000 pas.
	
	\textbf{Le solveur est validé. On peut allumer la barrière.}
	
	\begin{comprendre}
		Cette validation fixe notre barre d'erreur : environ 4\,\%.
		Toutes les mesures suivantes ont au moins cette incertitude.
		Il faudra s'en souvenir quand Mattéo comparera les temps tunnel.
	\end{comprendre}
	
	% ------------------------------------------------------------
	\clearpage
	\diapo{8}{Évolution face à la barrière (1 minute)}
	
	\textbf{Texte à dire :}
	
	Voici le film de la simulation, en quatre instants.
	
	Le paquet arrive de la gauche.
	Il heurte la barrière, la zone grisée.
	Il se sépare en deux :
	
	\begin{itemize}
		\item \textbf{un paquet réfléchi, dominant, qui repart vers la gauche ;}
		\item \textbf{un petit paquet transmis, qui continue vers la droite.}
	\end{itemize}
	
	Ce petit paquet transmis, c'est l'effet tunnel en images.
	\textbf{Classiquement, avec $E < V_0$, la particule serait toujours réfléchie.}
	Il n'y aurait rien du tout à droite.
	
	En bleu : le paquet libre, pour comparaison.
	Il avance à la même vitesse de groupe.
	C'est cette comparaison qui servira à mesurer l'avance ou le retard
	du paquet transmis.
	
	\begin{comprendre}
		Pendant l'impact, des franges apparaissent devant la barrière.
		C'est l'interférence entre le paquet qui arrive et le paquet qui repart.
		Bon détail à donner si le jury demande de commenter la figure.
	\end{comprendre}
	
	% ------------------------------------------------------------
	\clearpage
	\diapo{9}{Probabilités de réflexion et de transmission (1 minute)}
	
	\textbf{Texte à dire :}
	
	Pour mesurer, on intègre $|\psi|^2$ de chaque côté de la barrière,
	au cours du temps. C'est le code à droite :
	un masque sélectionne la zone après la barrière,
	une ligne intègre la densité dessus.
	
	La courbe du haut le confirme : \textbf{la norme totale reste constante.}
	C'est l'unitarité de Crank--Nicolson.
	
	À la fin, la probabilité d'être passé à droite --- c'est exactement ce que
	calcule cette ligne --- vaut $8{,}1\times10^{-2}$.
	La théorie des états stationnaires prédit $|T(k_0)|^2 = 2{,}5\times10^{-2}$.
	\textbf{Trois fois plus que la théorie. Est-ce une erreur ? Non.}
	C'est un de nos résultats les plus intéressants.
	
	L'explication :
	
	\begin{itemize}
		\item notre paquet n'est pas une onde plane ;
		\item il contient toute une gamme de vitesses autour de $k_0$ ;
		\item et la transmission $|T|^2$ monte très vite avec $k$.
	\end{itemize}
	
	\textbf{La barrière laisse surtout passer les composantes rapides.
		Elle agit comme un filtre.}
	
	Ce filtrage est la clé de toute la suite.
	C'est lui qui explique l'écart entre le temps mesuré
	et la théorie de Hartman, que Mattéo va vous montrer.
	
	Mattéo, à toi.
	
	\begin{comprendre}
		Image simple : la barrière reçoit un mélange de composantes lentes
		et rapides, et ne laisse passer presque que les rapides.
		
		Le paquet transmis est donc plus rapide, en moyenne, que le paquet
		incident. Et la probabilité mesurée dépasse la prédiction faite
		avec la seule vitesse moyenne $k_0$.
	\end{comprendre}
	
	% ------------------------------------------------------------	
\end{document}